% ============================================================
% LSE EC423 Extended Essay Template
% ============================================================
% 格式要求：A4，双倍行距，12pt，页码编号，2.5cm 页边距
% 字数上限 6,000（abstract、脚注、参考文献、附录不计入）
% ============================================================

\documentclass[a4paper, 12pt]{article}

% ----- 基础包 -----
\usepackage[utf8]{inputenc}
\usepackage[T1]{fontenc}
\usepackage{lmodern}              % Latin Modern（正文 + 数学）
\usepackage[british]{babel}     % British English 拼写规则
\usepackage{csquotes}

% ----- 页面布局 -----
\usepackage[
  a4paper,
  top=2.5cm, bottom=2.5cm,
  left=2.5cm, right=2.5cm
]{geometry}

\usepackage{setspace}
\doublespacing                  % 双倍行距

\usepackage{fancyhdr}
\pagestyle{fancy}
\fancyhf{}
\renewcommand{\headrulewidth}{0pt}
\fancyfoot[C]{\thepage}         % 页码居中

% ----- 数学 -----
\usepackage{amsmath, amssymb, amsthm}

% ----- 表格与图表 -----
\usepackage{booktabs}
\usepackage{graphicx}
\usepackage{float}
\usepackage{caption}
\captionsetup{font={small, singlespacing}, labelfont=bf}

% ----- 参考文献 -----
\usepackage[
  style=authoryear,
  backend=biber,
  natbib=true,
  maxcitenames=2,
  maxbibnames=99,
  uniquelist=false,
  giveninits=true,
  doi=false,
  isbn=false,
  url=false,
  eprint=false
]{biblatex}
\addbibresource{refs.bib}

% ----- 超链接（最后加载）-----
\usepackage[hidelinks]{hyperref}

% ----- 自定义命令 -----
\newcommand{\candidatenumber}{XXXXX}   % ← 替换为你的 5 位 candidate number
\newcommand{\coursecode}{EC423}

% ============================================================
\begin{document}

% ----- 封面页 -----
\begin{titlepage}
  \centering
  \vspace*{3cm}
  {\Large \textbf{\coursecode{} Extended Essay} \par}
  \vspace{1cm}
  {\large Independent Research \par}
  \vspace{2cm}
  {\LARGE \textbf{Canaries in Every Mine? \\ Offsetting Effects of AI on 
  Junior and Senior Workers within Occupations} \par}
  \vspace{2cm}
  {\large Candidate Number: \candidatenumber \par}
  \vspace{1cm}
  {\large MSc Econometrics and Mathematical Economics \par}
  {\large London School of Economics and Political Science \par}
  \vspace{1cm}
  {\large May 2026 \par}
  \vfill
  {\small Word count: [XXXX] \par}
\end{titlepage}

% ----- Abstract（不计入字数）-----
\newpage
\begin{abstract}
\noindent
[Your abstract here.]
\end{abstract}

\newpage
\tableofcontents
\newpage

% ============================================================
% 正文开始
% ============================================================

\section{Introduction}
\label{sec:intro}



\section{Theoretical Framework}
\label{sec:theory}

Consider a single occupation producing output $Y$, by two types of tasks:

\begin{itemize}
  \item Task $L$: Low expertise, including routine jobs, easy be replaced by AI automation.
  \item Task $H$: High expertise, including context-dependent tasks, management tasks, and other tasks requiring accumulated experience.

\end{itemize}

The output $Y$ follows a CES aggregation: 

$$
Y = [\beta y_L^\frac{\sigma-1}{\sigma} + (1-\beta)y_H^\frac{\sigma-1}{\sigma} ]^\frac{\sigma}{\sigma-1}
$$




\section{Existing Evidence}
\label{sec:literature}



\section{Data and Empirical Strategy}
\label{sec:empirics}

\subsection{Data: UK Labour Force Survey}

\subsection{AI Exposure Measurement}

\subsection{Empirical Specification}

\section{Results}
\label{sec:results}

\subsection{Descriptive Evidence}

\subsection{Main Results}

\subsection{Robustness}

\section{Discussion}
\label{sec:discussion}



\section{Conclusion}
\label{sec:conclusion}



% ============================================================
% 参考文献（不计入字数）
% ============================================================
\newpage
\singlespacing                  % 参考文献单倍行距，节省空间
\printbibliography

% ============================================================
% 附录（不计入字数）
% ============================================================
\newpage
\appendix
\section{Additional Tables and Figures}
\label{app:tables}



\end{document}